\emph{Översikt}
Den här övningen övar klasser och objekt på en och samma uppgift nästan
hela vägen: en klass som håller uppgifterna om en bok och hur långt du
har läst i den.  Vi börjar med att skriva klassen, lägger sedan till en
metod som inte får låta dig läsa längre än till sista sidan och följer
vad som händer med två objekt när bara det ena får ett anrop.  Därefter
kopplar vi in Pythons inbyggda operationer på våra egna objekt, en
dundermetod i taget: \mintinline{python}{<}, \mintinline{python}{==} och
sist de fyra räknesätten på en klass för bråk.  Då visar sig arvet: de
metoder du har skrivit ersatte metoder som klassen redan hade, och att
skriva en underklass är samma drag mot en klass du själv har skrivit.
Till sist väger vi en klass mot en uppslagslista, granskar ett program
som fungerar men är illa skrivet, och väljer --- som fördjupning ---
mellan arv och komposition.  Varje uppgift följs av ett
lösningsförslag --- ett av flera möjliga --- som du kan jämföra ditt
eget svar med.

\emph{Lärandemål}
Övningen övar de lärandemål som föreläsningen \emph{Klasser och objekt}
ställer upp, med två tillägg.  Målet om dundermetoder omfattar här
också \mintinline{python}{__eq__}, som föreläsningen beskriver men inte
implementerar och som modulens mål räknar upp, och dessutom de metoder
som ligger bakom räknesätten och typomvandlingen.  De senare behandlas
fullt ut i föreläsningen \emph{Operatoröverlagring}; här möter du dem i
en uppgift som låter dig leta upp dem i Pythons dokumentation, eftersom
mönstret är detsamma som du redan har använt.  Efter övningen ska
studenten kunna:
% Lo09 (sammansatta datatyper), Lo03 (dela upp program)
\begin{restatable}{lo}{OvningClassesLOAttributes}%
  \label{OvningClassesLOAttributes}%
  Skapa klasser med attribut och metoder.
\end{restatable}
% Lo09 (sammansatta datatyper)
\begin{restatable}{lo}{OvningClassesLOSelf}\label{OvningClassesLOSelf}%
  Förklara vad parametern \mintinline{python}{self} är, och använda den i
  metoder.
\end{restatable}
% Lo09 (sammansatta datatyper)
\begin{restatable}{lo}{OvningClassesLOObject}%
  \label{OvningClassesLOObject}%
  Förklara skillnaden mellan en klass och ett objekt.
\end{restatable}
% Lo09 (sammansatta datatyper)
\begin{restatable}{lo}{OvningClassesLOEncapsulation}%
  \label{OvningClassesLOEncapsulation}%
  Använda inkapsling för att skydda ett objekts data.
\end{restatable}
% Lo09 (sammansatta datatyper)
\begin{restatable}{lo}{OvningClassesLODunder}%
  \label{OvningClassesLODunder}%
  Implementera dundermetoder så att egna objekt kan skrivas ut,
  jämföras med varandra, räknas med och typomvandlas:
  \mintinline{python}{__init__}, \mintinline{python}{__str__},
  \mintinline{python}{__lt__}, \mintinline{python}{__eq__}, de fyra
  räknesätten och \mintinline{python}{__float__}.
\end{restatable}
% Lo01 (program utan kodupprepning), Lo09 (sammansatta datatyper)
\begin{restatable}{lo}{OvningClassesLOInheritance}%
  \label{OvningClassesLOInheritance}%
  Skapa underklasser med arv.
\end{restatable}
% Lo09 (sammansatta datatyper)
\begin{restatable}{lo}{OvningClassesLOOverride}%
  \label{OvningClassesLOOverride}%
  Överlagra metoder i underklasser, och anropa föräldraklassens metoder
  med \mintinline{python}{super()}.
\end{restatable}
% Lo09 (sammansatta datatyper), Lo06 (flexibla program)
\begin{restatable}{lo}{OvningClassesLORepresentation}%
  \label{OvningClassesLORepresentation}%
  Välja mellan att representera data med en klass eller med en
  uppslagslista, och motivera valet.
\end{restatable}
% Lo11 (granska program), Lo09 (sammansatta datatyper),
% Lo01 (program utan kodupprepning)
\begin{restatable}{lo}{OvningClassesLOReview}%
  \label{OvningClassesLOReview}%
  Granska en klass som andra har skrivit och föreslå förbättringar av
  namn, inkapsling och kodupprepning.
\end{restatable}
En sista, fördjupande uppgift övar dessutom:
% Lo09 (sammansatta datatyper)
\begin{restatable}{lo}{OvningClassesLOComposition}%
  \label{OvningClassesLOComposition}%
  Avgöra när arv respektive komposition passar bäst.
\end{restatable}

\emph{Förkunskaper}
Studenten bör ha tagit del av föreläsningen \emph{Klasser och objekt}:
den ger klass, objekt, attribut, metod, parametern
\mintinline{python}{self}, privata attribut, arv och dundermetoderna
\mintinline{python}{__init__}, \mintinline{python}{__str__} och
\mintinline{python}{__lt__}, som uppgifterna här bygger vidare på.
Därtill behövs föreläsningarna \emph{Funktioner}, \emph{Villkor och
styrstrukturer}, \emph{Upprepningar}, \emph{Moduler och paket},
\emph{Behållare: Listor} och \emph{Behållare: Uppslagslistor}:
studenten ska kunna skriva en funktion med parametrar och returvärde,
importera en egen modul, använda villkor och upprepningar och lagra
data i listor och uppslagslistor.  Särfall
(\mintinline{python}{raise} och \mintinline{python}{except}) förekommer
i tre lösningsförslag; de behandlas i föreläsningen \emph{Inmatning och
felhantering}.  Dundermetoderna bakom räknesätten och
\mintinline{python}{float} kräver inga förkunskaper utöver mönstret:
\cref{sec:brak} pekar på de avsnitt i Pythons dokumentation där namnen
står, och föreläsningen \emph{Operatoröverlagring} går igenom dem i
sin helhet.
